\paragraph{Reference geometry.}
We screen staging metadata for clear lesion links, a 3D centroid in each visit, and valid geometry before retrieving volumes; all 144 allocated volumes are retrieved. A separate sensitivity gives each patient equal weight. The 98 audited centroids span 12 patients. The median native axial spacing is 3 mm. Repeating the accuracy calculation in native earlier coordinates checks whether scaling from the shared affine changes the result.

\paragraph{Shared inputs for different matching mechanisms.}
The spline can optimize a continuous map on image samples that include the lesion, whereas the discrete search evaluates translations using outer anatomy and uses the lesion-bearing patches only in its scorer. SAM's own ranked list starts from the global maximum and adds separated locations in score order. Evaluating it separately measures the field's available coverage. For the shared-candidate experiment we take the geometry-ranked lists and sample SAM on exactly those centers. The two scorers therefore see the same opportunities at each budget.

\paragraph{Matched workloads.}
We time each method on the same 24 prepared staging queries and the same 32 prepared screening queries in separate cohort replays. Runs use a pinned CPU thread for the direct search and spline. The matched native replay measures eight screening pairs, four of which supply two queries. Native medians include the shared affine. The supplemental timing tables include stage times, p90 latency, and separate four-query estimates. A three-repeat randomized replay measures the smaller timing differences between fixed and adaptive search routes.
